% Generated draft: requires mathematical and source review.
\documentclass[11pt]{article}
\usepackage[margin=25mm]{geometry}
\usepackage{fontspec}
\setmainfont{Latin Modern Roman}
\usepackage{amsmath,amssymb,mathtools,unicode-math}
\setmathfont{Latin Modern Math}
\usepackage{enumitem,needspace,xurl,hyperref}
\setlength{\parindent}{0pt}
\setlength{\parskip}{5pt}
\setlength{\emergencystretch}{3em}
\newcommand{\N}{\mathbb{N}}
\newcommand{\Z}{\mathbb{Z}}
\newcommand{\Q}{\mathbb{Q}}
\newcommand{\R}{\mathbb{R}}
\newcommand{\C}{\mathbb{C}}
\newcommand{\F}{\mathbb{F}}
\newcommand{\A}{\mathbb{A}}
\newcommand{\PP}{\mathbb{P}}
\newcommand{\E}{\mathbb{E}}
\newcommand{\eps}{\varepsilon}
\newcommand{\dd}{\,\mathrm{d}}
\begin{document}
\section*{UnsolvedMath ID: 10400165}
\textbf{Source problem code:} catalogue entry
\subsection{Definitions and mathematical statement}
The source asks for three classifications.
\begin{enumerate}
\item Find and classify all semisimple monoidal categories having only finitely many isomorphism classes of simple objects. Here semisimplicity means decomposition into simple objects in the relevant category, and the monoidal structure consists of tensor product, unit and coherent associativity and unit constraints. The source does not specify the ground field, linearity, dimensions of morphism spaces, rigidity or simplicity of the unit; these hypotheses must not be silently added by replacing the class with a narrower class of fusion categories.
\item Find and classify finite-dimensional fusion-rule algebras and sets of \(6j\)-symbols. In the usual categorical setting, representatives \(X_i\) of the finitely many simple classes give tensor-product multiplicities through
\[
 X_i\otimes X_j\cong\bigoplus_k N_{ij}^{k}X_k,
\]
and \(6j\)-symbols express associativity maps in bases of the relevant multiplicity spaces, with coherence given by the pentagon equations. This describes the standard interpretation; the excerpt does not fix the coefficient field, admissible fusion rules, extra duality or positivity axioms, or equivalence of symbol systems under basis changes. It also does not state whether only realizable fusion rules are to be classified.
\item Find and classify all finite-depth subfactors. A subfactor is an inclusion \(N\subset M\) of von Neumann factors. In the standard finite-index setting, finite depth refers to finiteness of the simple sectors generated by the inclusion and its conjugate, equivalently finiteness of the corresponding principal-graph depth. The excerpt leaves the types of factors, index hypotheses and classification equivalence unspecified.
\end{enumerate}
These are requests for complete descriptions, with existence and equivalence criteria, rather than a conjectured explicit list. Monoidal equivalence in (1), changes of multiplicity-space bases in (2), and isomorphism of inclusions in (3) are natural comparison notions, but the exact notions intended require clarification.
\subsection{Short English statement}
Classify semisimple monoidal categories with finitely many simple classes, finite-dimensional fusion-rule algebras and \(6j\)-symbol systems, and finite-depth subfactors, keeping their respective equivalence conventions explicit.
\subsection{Sources}
\begin{itemize}
\item T. Ohtsuki, \emph{Problems on invariants of knots and 3-manifolds} (2002), Problem 9.1, supplied PDF page 135. \url{https://msp.org/gtm/2002/04/gtm-2002-04-024s.pdf}. Recorded access 2026-07-29; extraction pdf-font-structure. No individual proposer or additional bibliography entry is supplied for this item.
\item Local UnsolvedMath registry, \texttt{data/unsolvedmath/problems.json}, permanent ID \(10400165\), exact source code \texttt{catalogue entry}. The supplied summary calls the category-classification problem open; the registry research classification is OPEN-TRIAGE.
\end{itemize}
\end{document}
